\section{Experimental Setup}
\label{sec:experimental_setup}

This section describes the experimental design used to evaluate the proposed calibrated surrogate-assisted optimization framework. The experiments are organized around three questions: whether the global layout-parameter surrogate can provide useful candidate assessment on unseen layouts, whether online local calibration can reduce layout-specific prediction bias, and whether the calibrated surrogate-assisted strategy can reduce true FEA calls while maintaining comparable solution quality.

\subsection{Experimental Objectives}
\label{subsec:experimental_objectives}

The experiments are designed to examine the three key components of the proposed framework.

First, the global surrogate models are evaluated on unseen shear wall layouts. The feasibility surrogate is assessed in terms of both classification performance and conservative screening ability, while the reinforcement surrogate is assessed in terms of prediction accuracy and ranking quality. To verify the role of layout representation, three input settings are compared: design parameters only, design parameters with global layout statistics, and the proposed layout-parameter graph representation.

Second, the effectiveness of online local calibration is investigated. Since each optimization task is performed for a fixed layout and design condition, FEA-verified samples generated during optimization can provide layout-specific feedback. The calibration experiments therefore examine whether a small number of locally verified samples can improve feasibility probability calibration, reduce reinforcement prediction bias, and improve candidate ranking in the current layout.

Third, the calibrated surrogate models are embedded into optimization procedures. The optimization experiments compare full FEA-based optimization with several surrogate-assisted evaluation strategies. The comparison focuses on the number of true FEA calls, feasible-solution discovery, and final material cost. All compared methods use the same design space and true evaluator, and the final best solution is always determined by FEA-verified results.

\subsection{Dataset and Structural Evaluation}
\label{subsec:dataset_structural_evaluation}

The dataset is constructed from engineering shear wall layouts and a parametric structural evaluation workflow. Each sample corresponds to a layout, a design condition, and a section-material design scheme. The input includes layout information and design parameters, and the supervised labels include the comprehensive feasibility label $y_f$ and the reinforcement usage $m_s$.

The dataset contains 143 shear wall layouts designed by professional engineers. Following the design-condition grouping used in the previous layout-generation study, the layouts are divided into three groups according to seismic intensity and building height. Group7-H1 corresponds to low-demand cases with seismic intensity 7, PGA of 0.10~g, and building height $H \leq 50$~m. Group7-H2 corresponds to medium-demand cases with the same seismic intensity but building height $H>50$~m. Group8 corresponds to high-demand cases with seismic intensity 8 and PGA of 0.20~g. These groups reflect the combined influence of seismic demand and structural height on shear wall section design and material usage.

For each layout, 5000 candidate design schemes are sampled from a predefined parameter space. The sampled variables include the number of stories, wall thicknesses in different vertical zones, main and secondary beam dimensions, concrete strength grade, seismic intensity, site class, and seismic group. For parameters defined along the building height, basic engineering constraints are imposed during sampling. For example, upper-zone wall thickness is not allowed to be larger than lower-zone wall thickness.

A true structural evaluator $\mathcal{E}$ is used to generate labels for all sampled candidates. For each candidate, the evaluator automatically constructs the structural model, performs FEA in OpenSees, extracts global structural responses and member forces, and conducts code-based member design and reinforcement calculation. The comprehensive feasibility label $y_f$ is determined by analysis validity, global response limits, and member-level design checks. The reinforcement usage label $m_s$ is obtained from the final member design results.

The resulting dataset contains approximately 715,000 FEA-labeled samples. A layout-level split is adopted to evaluate cross-layout generalization. Specifically, 100 layouts are used for training, 21 layouts for validation, and 22 layouts for testing, corresponding to an approximate ratio of 0.70:0.15:0.15. All parameter variants of the same layout are assigned to the same subset to avoid information leakage between training and testing.

The parameter sampling space is summarized in Table~\ref{tab:parameter_space}. The detailed distribution of layouts and samples across design-condition groups is given in Table~\ref{tab:dataset_split}.

\begin{table}[t]
\centering
\caption{Sampling space of structural design parameters.}
\label{tab:parameter_space}
\footnotesize
\setlength{\tabcolsep}{3pt}
\begin{tabularx}{\columnwidth}{@{}>{\raggedright\arraybackslash}p{0.30\columnwidth}>{\centering\arraybackslash}p{0.17\columnwidth}>{\centering\arraybackslash}p{0.08\columnwidth}>{\raggedright\arraybackslash}X@{}}
\toprule
Parameter & Symbol & Unit & Sampling values \\
\midrule
Number of stories & $N$ & - & 18--33 \\
Story height & $h_{story}$ & m & 2.9 \\
Bottom wall thickness & $t_{w,b}$ & mm & 200, 250, 300, 350, 400 \\
Middle wall thickness & $t_{w,m}$ & mm & 160, 180, 200, 250, 300 \\
Top wall thickness & $t_{w,t}$ & mm & 160, 180, 200, 250 \\
Main beam depth & $h_{b,main}$ & mm & 400, 500, 550, 600, 650, 700 \\
Main beam width & $b_{b,main}$ & mm & 200, 250, 300, 350 \\
Secondary beam depth & $h_{b,sec}$ & mm & 300, 400, 450, 500 \\
Secondary beam width & $b_{b,sec}$ & mm & 200, 250, 300 \\
Concrete grade & $f_c$ & - & C30, C35, C40, C45, C50 \\
Seismic intensity & - & - & 6.0, 7.0, 7.5, 8.0 \\
Site class & - & - & I0, I, II, III, IV \\
Seismic group & - & - & 1, 2, 3 \\
\bottomrule
\end{tabularx}
\end{table}

\begin{table*}[t]
\centering
\caption{Dataset split and design-condition groups.}
\label{tab:dataset_split}
\begin{minipage}[c][0.16\textheight][c]{0.92\textwidth}
\centering
Placeholder for the dataset split table.\
Replace this box with the actual table.
\end{minipage}
\end{table*}

\subsection{Global Surrogate Model Training}
\label{subsec:global_surrogate_training}

Two global surrogate models are trained: the feasibility surrogate $SM_f$ and the reinforcement usage surrogate $SM_s$. Both models use the same candidate representation and network backbone, but they are trained for different prediction tasks. The feasibility surrogate outputs the probability that a candidate can pass the comprehensive structural checks, while the reinforcement surrogate predicts the reinforcement quantity before true evaluation.

To examine the effect of layout representation, three model settings are compared:
\begin{enumerate}
\item \textbf{Design-parameter MLP}: only section, material, and design-condition parameters are used as input;
\item \textbf{Layout-statistics MLP}: design parameters and global layout statistics are used as input;
\item \textbf{LayoutParamGNN}: room-level graph, global layout statistics, and design parameters are jointly used as input.
\end{enumerate}

The feasibility surrogate is trained using weighted binary cross-entropy loss to account for class imbalance. The validation set is used to select the screening threshold under a target feasible recall, rather than using the default threshold of 0.5. This setting is consistent with the role of the feasibility surrogate as a conservative pre-screening tool.

The reinforcement surrogate is trained using the logarithmically transformed and standardized reinforcement quantity as the regression target. Smooth L1 loss is used to reduce the influence of large outliers. In addition to ordinary regression metrics, ranking metrics are also reported because the reinforcement surrogate is mainly used to rank candidate designs by material economy during optimization.

The main training configurations and hyperparameters are summarized in Table~\ref{tab:surrogate_training_settings}.

\begin{table*}[t]
\centering
\caption{Training settings of global surrogate models.}
\label{tab:surrogate_training_settings}
\begin{minipage}[c][0.16\textheight][c]{0.92\textwidth}
\centering
Placeholder for surrogate model training settings.\
Replace this box with the actual table.
\end{minipage}
\end{table*}

\subsection{Local Calibration Settings}
\label{subsec:local_calibration_settings}

Local calibration experiments are conducted on the test layouts to simulate the use of FEA-verified samples accumulated during optimization. For each test layout, a certain number of verified samples are randomly selected as the local calibration set, and the remaining samples are used as the layout-specific evaluation set. Each local calibration sample stores the candidate representation, the global surrogate predictions, and the true FEA-based labels.

For feasibility calibration, the input is the global feasibility prediction, or its logit transformation, and the target is the true feasibility label. Several lightweight calibration methods can be compared, including Platt scaling, isotonic regression, local logistic calibration, and probability residual correction. The calibrated probabilities are evaluated using probability calibration metrics and conservative screening metrics.

For reinforcement calibration, the target is the residual between the true reinforcement quantity and the global surrogate prediction. Candidate features and the global reinforcement prediction are used to fit a local residual model. Candidate methods include ridge regression, Gaussian process regression, K-nearest neighbors, gradient boosting, and a small multilayer perceptron. The calibrated reinforcement predictions are evaluated by regression accuracy, systematic bias, and ranking quality.

To examine sample efficiency, different local calibration sizes are considered, such as 25, 50, 100, 200, and 500 samples. For each size, repeated random sampling is conducted, and average results are reported. The local calibration settings are summarized in Table~\ref{tab:local_calibration_settings}.

\begin{table*}[t]
\centering
\caption{Settings for online local calibration experiments.}
\label{tab:local_calibration_settings}
\begin{minipage}[c][0.16\textheight][c]{0.92\textwidth}
\centering
Placeholder for local calibration settings.\
Replace this box with the actual table.
\end{minipage}
\end{table*}

\subsection{Optimization Benchmark Settings}
\label{subsec:optimization_benchmark_settings}

Optimization experiments are performed on the test layouts to evaluate whether the proposed framework can reduce true FEA calls in actual search processes. For each test layout, the layout $L$ and design condition $c$ are fixed, and the optimization variables are the section and material parameters $x$. All compared methods use the same design space, true evaluator, and random seeds when applicable.

The main optimizer is a genetic algorithm. To test whether the surrogate-assisted mechanism depends on a specific optimizer, particle swarm optimization and random search are also included as supplementary optimizers. For all optimizers, the final feasibility and material cost of each selected design are obtained from the true evaluator $\mathcal{E}$.

Four candidate evaluation strategies are compared:
\begin{enumerate}
\item \textbf{Full FEA optimization}: all candidates generated by the optimizer are evaluated by $\mathcal{E}$. This strategy serves as the baseline.
\item \textbf{Cost surrogate}: candidates are ranked using the reinforcement surrogate and estimated concrete quantity, and only high-priority candidates are submitted to $\mathcal{E}$.
\item \textbf{Screen + cost surrogate}: feasibility screening is applied before cost-based ranking.
\item \textbf{Calibrated screen + cost surrogate}: online local calibration is further introduced, and calibrated feasibility and reinforcement predictions are used for screening and ranking.
\end{enumerate}

For surrogate-assisted strategies, each optimizer first generates a candidate batch $\mathcal{B}_t$. The surrogate module then evaluates all candidates in the batch and selects a subset $\mathcal{S}_t$ for true FEA and code checking. The verified results are used to update the optimizer and the local calibration set. Candidates not selected for FEA are not used as true samples and are not considered in final best-solution comparison.

The optimization benchmark settings are summarized in Table~\ref{tab:optimization_settings}.

\begin{table*}[t]
\centering
\caption{Optimization benchmark settings and compared candidate evaluation strategies.}
\label{tab:optimization_settings}
\begin{minipage}[c][0.18\textheight][c]{0.92\textwidth}
\centering
Placeholder for optimization benchmark settings.\
Replace this box with the actual table.
\end{minipage}
\end{table*}

\subsection{Evaluation Metrics}
\label{subsec:evaluation_metrics}

The evaluation metrics are divided into three groups: global surrogate metrics, local calibration metrics, and optimization performance metrics.

For the feasibility surrogate, ROC-AUC, PR-AUC, and F1 score are used to evaluate overall classification performance. Since the model is used as a conservative screening tool, feasible recall, reject rate, and false reject count are also reported. Feasible recall measures the proportion of truly feasible candidates retained after screening. Reject rate measures the proportion of candidates filtered before FEA. False reject count measures the number of feasible candidates mistakenly filtered out.

For the reinforcement surrogate, mean absolute error (MAE), root mean square error (RMSE), coefficient of determination ($R^2$), mean absolute percentage error (MAPE), and bias are used to evaluate regression performance. Since reinforcement prediction is mainly used for candidate ranking, pairwise ranking accuracy, Spearman correlation, and top-$k$ recall are additionally used to evaluate ranking quality.

For local calibration, feasibility probability calibration is evaluated using Brier score and expected calibration error (ECE), together with screening recall and reject rate. Reinforcement residual calibration is evaluated using MAE, RMSE, bias, and ranking metrics after calibration.

For optimization experiments, computational efficiency is measured by the number of true FEA calls and the FEA reduction ratio. Solution quality is measured by feasible-solution discovery rate, the number of FEA calls required to find the first feasible solution, and final material cost. For each test layout and random seed, surrogate-assisted strategies are compared with full FEA optimization to determine whether FEA reduction is achieved without substantial degradation in feasibility discovery or final cost.

The evaluation metrics are summarized in Table~\ref{tab:evaluation_metrics}.

\begin{table*}[t]
\centering
\caption{Evaluation metrics for global surrogate models, local calibration, and optimization performance.}
\label{tab:evaluation_metrics}
\begin{minipage}[c][0.18\textheight][c]{0.92\textwidth}
\centering
Placeholder for evaluation metrics.\
Replace this box with the actual table.
\end{minipage}
\end{table*}
